\chapter{Classical and Quantum Coding Theory}

\section{Classical Codes}

In classical information theory, a code is defined as a subset $M \subseteq \mathbb{B}^n$ containing $2^m$ elements (\cite{Kitaev2002}). Notably, $m$ is not restricted to integer values; for a code with six possible values, $m = \log_2(6) \approx 2.5849$. The difference between the physical block length $n$ and the message length $m$ defines the redundancy of the system. This redundancy is a fundamental requirement for verifying the integrity of a received bitstring. While these subsets typically map to external data structures, such as the ASCII standard where $m = 7$, they more broadly function as the essential interface between physical bit states and logical information.

In the context of error correction, codes are structured such that the underlying information remains recoverable despite noise. Classical errors are typically discrete, manifesting as ``bit-flips'' within a string. Consequently, techniques such as data redundancy and the cloning of information are central to classical error mitigation. It is important to note that these specific strategies are prohibited in the quantum regime due to the No-Cloning Theorem.

\section{Quantum Codes}

In quantum computation, defining a code as a discrete subset of bitstrings is insufficient, as it fails to account for the continuous nature of quantum states. Instead, a quantum code is defined as a subspace of the total Hilbert space (\cite{Kitaev2002}). This formulation allows for the encoding of arbitrary superpositions of basis states. By leveraging the linearity of vector spaces, one can define an encoding transformation by mapping the original basis vectors into a linearly independent set that spans the protected code subspace.

This subspace-based approach ensures that linear transformations can be applied while the state remains within the code's boundaries. Furthermore, it provides the mathematical framework necessary to identify and rectify errors without collapsing the logical state. This is achieved by extracting an error syndrome, a diagnostic indicator of the noise process, via projective measurements onto the subspace, rather than by measuring individual physical qubits.

\section{Recovery Conditions}

To maintain generality in the theory of quantum error-correcting codes, minimal assumptions are made regarding the noise model. We consider an error channel $\mathcal{E}$ acting on a density matrix $\rho$ as a trace-preserving map:

\begin{equation}
\mathcal{E}(\rho) = \sum_{i} E_i \rho E_i^\dagger
\end{equation}

where $\{E_i\}$ are the Kraus operators representing the noise process. A code $C$ is said to correct the channel $\mathcal{E}$ if there exists a recovery operation $\mathcal{R}$ such that (\cite{Nielsen2010}):

\begin{equation}
\forall \rho \in C \implies (\mathcal{R} \circ \mathcal{E})(\rho) \propto \rho
\end{equation}

The map $\mathcal{R}$ is a physical process designed to map the corrupted state back into the code space. Critically, the recovery is considered successful if the logical information is preserved, even if the resulting physical state is not identical to the original codeword. This distinction is vital for understanding degenerate codes, where multiple distinct error processes may lead to the same corrected state within the logical subspace.

\section{Knill-Laflamme Conditions}

The existence of a recovery operation for an error is not guaranteed for an arbitrary code. Therefore, it is crucial to have a method to determine if a code can correct an error $\mathcal{E}$ within the subspace spanned by $\{E_i\}$. For a code $\mathcal{C}$ with projector $P$, the necessary and sufficient condition for the code to correct the error $\mathcal{E}$ is given by the Knill-Laflamme conditions (\cite{Nielsen2010}):

\begin{equation}
P E^\dagger_i E_j P = \alpha_{ij} P
\end{equation}

for some Hermitian matrix $\alpha$.

\textbf{Proof (adapted from Nielsen and Chuang \citep{Nielsen2010}):}

To prove sufficiency, we can construct the operation $\mathcal{R}$ assuming the Knill-Laflamme conditions are fulfilled. We assume a general method for error correction that separates syndrome measurement and error recovery, proving that both components can be generated given these conditions.

By assumption, $\alpha$ is a Hermitian matrix and is therefore diagonalizable as $d = u^\dagger \alpha u$. We can define operations $F_k := \sum_{i} u_{ik} E_i$, and given Theorem~\ref{theorem_8_2}, $\{F_k\}$ is also a set of operation elements for $\mathcal{E}$. By substitution, we get:

\begin{align*}
P F^{\dagger}_k F_l P &= \sum_{ij} u_{ki}^\dagger u_{jl} P E_i^\dagger E_j P\\
&= \sum_{ij} u_{ki}^\dagger u_{jl} \alpha_{ij} P\\
&= \sum_{ij} u_{ki}^\dagger \alpha_{ij} u_{jl} P\\
&= d_{kl} P
\end{align*}

Using the polar decomposition for $F_k P$, we obtain:

\begin{align*}
F_k P &= U_k\sqrt{(F_k P)^\dagger F_k P}\\
&= U_k\sqrt{P F_k^\dagger F_k P}\\
&= U_k\sqrt{d_{kk} P}\\
&= \sqrt{d_{kk}} U_k P
\end{align*}

Therefore, $F_k$ rotates into a subspace projected by $P_k := U_k P U_k^\dagger = \frac{F_k P U_k^\dagger}{\sqrt{d_{kk}}}$. Moreover, given that $d$ is a diagonal matrix, it is evident that the spaces projected by $P_k$ are orthogonal, allowing for a reliable syndrome measurement. This could also be augmented into a trace-preserving operation by adding the necessary subspaces (though this is not crucial in every code, such as the repetition code). The error correction step would simply involve applying $U^\dagger_k$.

Thus, we can construct the operation $\mathcal{R}(\sigma) = \sum_k U_k^\dagger P_k \sigma P_k U_k$. We can show that it satisfies the recovery condition because, without loss of generality for $\rho \in \mathcal{C}$:

\begin{align*}
\mathcal{R}(\mathcal{E}(\rho)) &= \sum_{kl} U_k^\dagger P_k F_l \rho F_l^\dagger P_k U_k\\
&=\sum_{kl} \left(U_k^\dagger P_k F_l \sqrt{\rho}\right) \left(\sqrt{\rho} F_l^\dagger P_k U_k\right)\\
&=\sum_{kl} \left(U_k^\dagger P_k F_l \sqrt{\rho}\right) \left(U_k^\dagger P_k F_l \sqrt{\rho}\right)^\dagger
\end{align*}

Evaluating the inner term:

\begin{align*}
U_k^\dagger P_k F_l \sqrt{\rho} &= U_k^\dagger P_k F_l P \sqrt{\rho}\\
&= \frac{U_k^\dagger U_k P F_k^\dagger F_l P \sqrt{\rho}}{\sqrt{d_{kk}}}\\
&= \frac{\delta_{kl} d_{kk} P \sqrt{\rho}}{\sqrt{d_{kk}}}\\
&= \delta_{kl} \sqrt{d_{kk}} \sqrt{\rho}
\end{align*}

Substituting this back yields:

\begin{align*}
\mathcal{R}(\mathcal{E}(\rho)) &=\sum_{kl} \left(\delta_{kl} \sqrt{d_{kk}} \sqrt{\rho}\right) \left(\delta_{kl} \sqrt{d_{kk}} \sqrt{\rho}\right)^\dagger\\
&=\sum_{kl} \delta_{kl} d_{kk} \rho\\
&\propto \rho
\end{align*}

This satisfies the conditions for $C$ to correct $\mathcal{E}$.

To prove the necessity of the Knill-Laflamme conditions, suppose $\mathcal{E}$ is perfectly correctable by the recovery operation $\mathcal{R}$ with operation elements $\{R_j\}$. Define an operation $\mathcal{E}_C (\rho) = \mathcal{E}(P\rho P)$ which is also recoverable by $\mathcal{R}$, where the proportionality factor is a constant $c$ independent of $\rho$. This can be rewritten as:

\begin{align*}
\mathcal{R}(\mathcal{E}_C(\rho)) &= c P\rho P\\
\sum_{ij} R_j E_i P \rho P E^\dagger_i R^\dagger_j &= c P \rho P
\end{align*}

By Theorem~\ref{theorem_8_2}, we know there exist complex numbers in a unitary matrix $c$ such that:

\begin{align*}
R_k E_i P &= c_{ki} P\\
P E_i^\dagger R_k^\dagger &= c_{ki}^* P\\
P E_i^\dagger R_k^\dagger R_k E_j P &= c_{ki}^* c_{kj} P
\end{align*}

Considering that $\mathcal{R}$ is trace-preserving, summing over $k$ yields:

\begin{equation}
P E_i^\dagger E_j P = \sum_k c_{ki}^* c_{kj} P = \alpha_{ij} P
\end{equation}

This perfectly matches the Knill-Laflamme conditions.

\section{Characteristics Imposed on a Code by the Knill-Laflamme Conditions}

The Knill-Laflamme conditions are a great way to reveal that, due to the existence of $\alpha$ as a Hermitian matrix and the conditions themselves, a quantum code that corrects an error needs to be:

\begin{enumerate}
\item \textbf{Orthogonal:} The basis should be orthogonal so that the states are distinguishable, which is reflected by the syndrome measurement operation obtained in the proof of the Knill-Laflamme conditions. This is also indispensable for $\alpha$ to be a Hermitian matrix (\cite{Kitaev2002}).
\item \textbf{Non-deformed by errors:} This corresponds, by the definition of deformation, to the Knill-Laflamme conditions. It basically represents that a codeword should not be changed or made non-orthogonal by the application of the error. Thus, even when the error is applied, it is still possible to distinguish between states and recover them.
\end{enumerate}

\begin{lemma}[Theorem 8.2, Nielsen and Chuang \citep{Nielsen2010}]\label{theorem_8_2}
Suppose $\{E_1, E_2, \ldots, E_m\}$ and $\{F_1, F_2, \ldots, F_n\}$ are operation elements giving rise to quantum operations $\mathcal{E}$ and $\mathcal{F}$ respectively. By appending zero operators to the shorter list, we may ensure that $m = n$. Then $\mathcal{E} = \mathcal{F}$ if and only if there exist complex numbers $u_{ij}$ such that $E_i = \sum_{j} u_{ij} F_j$, where $u$ is an $m \times m$ unitary matrix.
\end{lemma}
